\chapter{Conclusion and Future Work}\label{ch:conclusion}

\section{What the work set out to do, and what it found}

This project began as the construction of a secure ESP32 intrusion alarm. It became a study of
why the techniques normally relied upon to establish that such a system is secure did not
establish it.

Four defects were found \emph{in this project's own artefacts}. Each had survived the instrument
that existed to catch it:

\begin{itemize}[leftmargin=1.6em,itemsep=4pt]
  \item A compile-time guard accepted transport security plus a broker login as evidence of
        publisher authority. It survived code review and emitted no diagnostic. The ablation
        found it, and measured its cost: alarm activation of 16.7\% against 100\% for the
        control, with the adversary disarming the device in 246 of 300 trials.
  \item A formal invariant stating that unauthenticated remote disarm never takes effect was
        \emph{unfalsifiable by construction}. It passed continuous integration on every push.
        It was incapable of failing.
  \item The safety loop had no static bound on its own execution time. The defect had passed
        9,600 adversarial trials, 3,000,000 fuzzing iterations, exhaustive model checking and 129
        host tests. The response-time decomposition found it, because the decomposition required
        a bound that did not exist.
  \item Four further defects were found in the measurement infrastructure itself, two of which
        produced results that flattered the intended conclusion.
\end{itemize}

The pattern common to all of them is a single sentence, and it is the work's central claim:

\begin{quote}
A security property asserted in build configuration, and a timing property asserted in design,
are both unverified claims. Each can be falsified by a method that costs far less than the system
it checks.
\end{quote}

The guard asserted authorisation in a preprocessor expression. The specification asserted a
refusal in a comment. The loop asserted a bound nowhere at all. In each case the assertion was
accepted because it appeared in a location where assertions are not normally examined.

\section{Objectives against outcomes}

\begin{longtable}{@{}p{0.05\textwidth}p{0.49\textwidth}p{0.39\textwidth}@{}}
\toprule
\textbf{} & \textbf{Objective} & \textbf{Outcome} \\
\midrule
\endhead
O1 & Adequacy of core pinning with bounded queue discipline, and what can be settled without
hardware & \textbf{Partially met.} Logical isolation proved exhaustively (660 states). Physical
timing invariance not measured, and the boundary between the two is now stated precisely
wherever the claim appears. \\
O2 & Formally, whether a remote adversary can affect the safety function's control or data flow
& \textbf{Met.} \code{SafetyCanAlwaysStep} holds in every reachable state under an
unfair, unboundedly stalling, crashing network task. \\
O3 & Which command-path defence preserves the safety function, availability and replay
resistance & \textbf{Met.} The three properties fail independently. Only the counter-bound
message authentication code holds all three; the rate limiter costs 41\% of operator commands.
\\
O4 & Where broker authorisation and application-layer authentication fail, and whether either
dominates & \textbf{Met, with a negative result.} Neither dominates. The access control list
contains a device compromise the signature cannot; the signature contains a controller
compromise the list cannot. The layered build falls to E$'$. \\
O5 & A response-time bound, separating analytical from hardware terms & \textbf{Partially met.}
Three of four interference terms removed by construction; the fourth requires hardware. The
derivation exposed the unbounded drain. \\
O6 & Whether the configuration defect generalises to comparable open-source projects &
\textbf{Attempted; null result.} Rule 1: zero hits across 135 projects. Rule 2: nine confirmed
cases, reported as an existence claim and not a prevalence estimate. \\
O7 & A falsifiable measurement protocol for the unresolved timing claim, published whether or not
it can be executed & \textbf{Met.} Pre-registered before any data existed, with the test,
sample sizes and acceptance criteria fixed in advance (Appendix~\ref{app:hardware}). \\
\bottomrule
\end{longtable}

Six of seven objectives are met or partially met. The two partial outcomes share one cause: the
absence of physical hardware, which is the work's principal limitation
(Chapter~\ref{ch:limitations}).

\section{What this work establishes}

\begin{enumerate}[leftmargin=2.0em,itemsep=4pt]
  \item \textbf{Security properties of a command path fail independently, and must be measured
        separately.} Alarm activation, operator availability and replay resistance were each
        preserved by a design that destroyed one of the others. A single-metric evaluation of
        these six designs would have reached the wrong conclusion, and the figure reporting them
        needs three panels for exactly that reason.
  \item \textbf{Transport authentication is not publisher authorisation in MQTT 3.1.1.} The
        broker re-originates messages; the protocol carries no publisher identity to the
        subscriber. The cost of the unsound inference was measured rather than argued.
  \item \textbf{Broker authorisation and application-layer authentication are complementary and
        neither subsumes the other.} Established against real broker software, because scenarios
        E and E$'$ carry the same valid signature and differ only in the transmitting identity
        --- a distinction no amount of reasoning about the cryptography would have settled.
  \item \textbf{A formal property must be shown to be falsifiable before its success is
        evidence.} The repair is not a better invariant; it is a configuration whose
        \emph{required} outcome is a violation, with the specific invariant name required in the
        output. The same reasoning produced the negative control for trace conformance.
  \item \textbf{Response-time analysis is a defect-finding instrument and not only a
        certification instrument.} Four other methods had passed the firmware; deriving the bound
        found the defect, because the derivation required a quantity the code did not have.
  \item \textbf{Measurement infrastructure requires the same scrutiny as the artefact.} Four
        defects were found in the harnesses, two of them flattering. Every repair is now
        accompanied by a control that would fail if the defect recurred.
\end{enumerate}

\section{What this work does not establish}

Restated compactly, because the conclusions should not be readable without it.

No timing measurement of any kind was taken. Physical timing invariance is unmeasured and
unclaimed; logical isolation does not imply it. On-device memory safety is not established, the
fuzzing having run on a host build under sanitizers that do not exist for the Xtensa target. No
radio-layer experiment was performed. The corpus study makes no prevalence claim. The flash-dump
containment depends on the access control list and on the device identity lacking command write
access, and not on the key being protected --- it is not. The layered configuration falls to a
compromised controller host. Two design problems identified by this work remain open and
unfixed: owner lockout under key compromise, and scenario E$'$.

\section{Future work}

\subsection{Immediate: execute the pre-registered hardware protocol}

Conditions A--D, as specified in Appendix~\ref{app:hardware} and
\code{docs/HARDWARE\_A\_TO\_C\_RUNBOOK.md}. The analysis plan is already fixed, so the measurement
can confirm or falsify the timing claim without the test being chosen after the data is seen. This
is the single highest-value next step, because it is the only one that resolves a claim the work
currently has to leave open. Conditions A--C require an ESP32 DevKit, a logic analyser or a second
microcontroller as a timing reference, and a controlled network; condition D additionally requires
an RF-shielded enclosure.

\subsection{Repair the two open design problems}

\textbf{Owner lockout.} A counter-resynchronisation path authenticated by the local button, which
is physically present and outside a remote adversary's reach. The obvious alternatives --- a
resynchronisation command over MQTT, or a tolerance window on the counter --- weaken the replay
resistance the counter exists to provide, so the design question is not trivial and the proposed
repair needs evaluation against the same twelve-arm ablation before it can be claimed to work.

\textbf{Scenario E$'$.} Mitigation requires an authority that does not reside on the controller
host: a hardware security module, a second factor for disarm, or a physically present
confirmation for state-changing commands. Each shifts the cost of the attack rather than
eliminating it, and the trade-off against operator availability --- already shown to be a real
cost --- would have to be measured.

\subsection{Extend the evaluation}

Flash encryption and secure boot, which would address key extraction rather than containing it.
Other broker implementations with different authorisation models --- EMQX, HiveMQ, AWS IoT Core
--- since scenario E's containment is contingent on Mosquitto's deny-by-default list. On-device
fuzzing, which requires an instrumentation approach that does not depend on host sanitizers. A
soak test for flash wear from counter commits, clock drift and fragmentation over months.

\subsection{Generalise the method rather than the finding}

The corpus study attempted to generalise the \emph{defect} and returned a null result. A more
promising direction is to generalise the \emph{instrument}: a static check for the class of
unverified configuration assumption found here, validated against a labelled set of positive and
negative controls rather than against a star-filtered GitHub sample. The eight controls built for
this project's rule, including its own pre-repair guard, are the starting point for such a
labelled set.

\section{Closing remark}

The most useful outcome of this project is not the firmware. It is the record of four methods
that each reported success while a defect they were meant to catch was present, and of the
specific, testable changes that make each method capable of failing. A test that cannot fail is
not evidence, and that applies to a build guard, a model-checked invariant, a fuzzing campaign and
a measurement harness alike.
